\begin{figure}[htbp]
\centering
\begin{tikzpicture}[x=.88cm,y=.88cm,font=\small]
  \fill[paper] (0,0) rectangle (12.4,6.25);
  \draw[source!75,line width=1.1pt] (.7,2.2) .. controls (4.2,1.3) and
    (8.1,5.2) .. (11.8,4.65);
  \node[text=source,below,font=\scriptsize] at (6.35,2.55) {长江（示意）};
  \fill[timeline!30] (1.65,3.95) ellipse (1.15 and .72);
  \node[align=center] at (1.65,3.95) {成都\\益州腹地};
  \fill[timeline!38] (4.45,3.42) ellipse (1.15 and .72);
  \node[align=center] at (4.45,3.42) {白帝\\退守节点};
  \fill[dynasty!30] (7.12,3.83) ellipse (1.24 and .78);
  \node[align=center] at (7.12,3.83) {夷陵、猇亭\\江汉关口};
  \fill[source!30] (10.55,3.95) ellipse (1.2 and .72);
  \node[align=center] at (10.55,3.95) {武昌\\吴军上游支点};
  \fill[source!22] (9.72,1.36) ellipse (1.18 and .68);
  \node[align=center,font=\scriptsize] at (9.72,1.36) {建业\\下游中枢};
  \draw[->,very thick,color=timeline] (2.66,3.83) -- (5.88,3.8);
  \node[text=timeline,above,font=\scriptsize] at (4.3,4.16) {蜀军东下（概括）};
  \draw[->,very thick,color=dynasty] (9.35,3.98) -- (8.38,3.91);
  \node[text=dynasty,above,font=\scriptsize] at (8.88,4.26) {吴军据水陆节点};
  \draw[->,thick,dashed,color=timeline] (6.2,3.35) .. controls (5.45,2.65)
    and (4.98,2.45) .. (4.63,2.75);
  \node[text=timeline,below,font=\scriptsize] at (5.85,2.48) {蜀军退向白帝};
  \draw[->,thick,dashed,color=source] (10.15,3.27) .. controls (10.03,2.57)
    and (9.93,2.12) .. (9.81,1.95);
  \node[text=source,right,font=\scriptsize] at (10.2,2.48) {上、下游调度};
  \node[draw=source!75,fill=paper,rounded corners=2pt,align=center,
    inner sep=3pt,font=\scriptsize] at (6.35,.65)
    {夷陵的胜负不能脱离江面、营地移动和补给线；箭线不复原战场细部。};
\end{tikzpicture}
\caption{夷陵、白帝与孙吴上游防务（222 年，示意）。据《三国志·先主传》、
《吴主传》及相关列传概括。地点的相对关系用于说明蜀军深入和吴军依托江汉
节点的选择；不表示营垒长度、火攻过程、伤亡或精确行军线。}
\label{fig:vol02b:yiling-river-defense}
\end{figure}
